\documentclass{tietreport}

% Import images
\usepackage{graphicx}
\usepackage{float}

% Bibliography configuration
\usepackage[
  date=year,
  style=alphabetic,
  backend=biber,
  sorting=nty,
  maxnames=5,
  minnames=3,
  maxalphanames=4,
  minalphanames=3,
  backref=true,
  doi=false,
  isbn=false,
  url=false,
  eprint=false
]{biblatex}
\addbibresource{references.bib}

%% ==================================================
%% CUSTOMIZE THESE FIELDS FOR YOUR PROJECT
%% ==================================================

% Course/Lab header
\TitlePageHeader{%
  UCS503P - Software Engineering Project%
}

% Main project title
\ProjectTitle{%
  Cardio AI: Explainable Heart Disease Prediction%
}

% Subtitle or project description
\ProjectSubTitle{%
  Using CNN-LSTM Attention Network with Grad-CAM Interpretability%
}

% Student information (up to 4 students)
\author{%
  \texttt{1024170434} Dhruv Singla \\
  \texttt{1024170451} Srishti%
}

% Group number, degree, year, program
\TitlePageSubText{%
  Group: BE Third Year, CSE%
}

% Faculty advisor/guide name
\AdvisorName{%
  Nisha Thakur%
}

% Evaluation period and department info
\TitlePageFooterText{{%
    \large Mid-Semester Evaluation \\[0.2cm]
    \bfseries Computer Science and Engineering
    Department%
  } \\[0.15cm]
  Thapar Institute of Engineering and Technology,
  Patiala%
}

%% ==================================================

\begin{document}

% Generate title page
\maketitle

% Table of contents (auto-generated from chapters)
\tableofcontents

% ===== CHAPTER 1: INTRODUCTION =====

\chapter{Introduction}

\section{Background and Context}

Cardiovascular disease remains the leading cause of mortality globally, necessitating advanced, early-prediction systems. While deep learning models can predict heart disease from routine clinical measurements---such as blood pressure, cholesterol levels, and ECG markers---their inherent "black box" nature poses a significant barrier to clinical adoption. Regulatory bodies and medical ethics demand interpretable, auditable predictions. This project, \textbf{Cardio AI}, aims to bridge this gap by prioritizing clinical transparency. Our objective is to translate complex deep neural network confidence scores into actionable, human-readable clinical narratives, empowering physicians to trust AI-assisted diagnoses.

\subsection{Problem Statement}

While standard classification tasks can achieve high accuracy on clinical datasets, the primary challenge addressed by this project is shifting the focus from simple heart disease classification to \textbf{clinical transparency}---making deep learning explainable rather than a black box. Current challenges in AI-driven heart disease prediction include:
\begin{itemize}
    \item \textbf{Black-Box Predictions:} Deep learning models, while highly accurate, offer no clinical reasoning or transparent decision pathways, leading to physician reluctance.
    \item \textbf{Feature Opacity:} The cardiovascular domain involves heterogeneous clinical features (e.g., numerical vitals, categorical ECG types, binary indicators). Understanding which specific physiological parameters drive a patient's risk score is non-trivial.
    \item \textbf{Operational Constraints \& Overfitting Risks:} Many theoretical models fail in practice due to overfitting on limited datasets (such as the UCI Cleveland dataset). Robust systems must deliver immediate value and incorporate strategies to fail fast through rapid iteration.
    \item \textbf{Lack of Visual Evidence:} Clinicians lack interactive visual artifacts (such as heatmaps or feature importance charts) that map model confidence directly back to patient physiology.
    \item \textbf{Absence of Longitudinal Tracking:} Existing diagnostic tools rarely support patient-history-aware predictions or longitudinal trend analysis across repeated assessments.
\end{itemize}

\section{Project Objectives}

To address these challenges, we have defined the following core objectives:
\begin{enumerate}
  \item \textbf{Develop an Explainable Deep Learning Model:} Architect and train a hybrid CNN-LSTM-Attention model on the UCI Cleveland Heart Disease dataset. This model will incorporate multiple Explainable AI (XAI) methods, specifically Grad-CAM, Attention weights, and SHAP (Shapley Additive exPlanations), to decode the model's reasoning.
  \item \textbf{Deliver an Interactive Clinical Dashboard:} Create a responsive, premium Single Page Application (SPA) using HTML5, CSS3, Bootstrap 5, and JavaScript. This dashboard will allow clinicians to input physiological parameters and visualize predictions via Plotly.js interactive charts.
  \item \textbf{Automate Clinical Reporting:} Generate automated PDF reports featuring AI-driven narrative summaries, feature-by-feature clinical interpretations, and guideline-referenced evaluations.
  \item \textbf{Implement Patient Management Architecture:} Design a robust backend using Flask (Python) and a MySQL database to securely track longitudinal patient records.
\end{enumerate}


% ===== CHAPTER 2: METHODOLOGY & DESIGN =====

\chapter{Methodology and Design}

\section{System Architecture}

Our proposed solution combines deep learning engineering with clinical interpretability, focusing on robust pipeline design and multi-method XAI.

\subsection{High-Level Design and Modeling}

During the initial phase of the project, significant effort was dedicated to defining actor modeling, system flow, and architectural blueprints.
\begin{itemize}
    \item \textbf{Deep Learning Pipeline:} The core prediction engine is designed around a CNN-LSTM-Attention network. Clinical data will be transformed into spatiotemporal formats using Gramian Angular Summation Field (GASF) image encoding for robust feature extraction.
    \item \textbf{Multi-Method Explainability Engine:} The architecture maps parallel processing flows to compute SHAP values, Grad-CAM 2D heatmaps, and Attention weights simultaneously after the CNN-LSTM inference.
    \item \textbf{UML Modeling:} We designed a comprehensive UML Use Case Diagram establishing boundaries and main features for the primary actor (Clinician) and secondary actors (Patient, ML Engine, MySQL Database). Features include predicting risk, viewing XAI charts, performing Exploratory Data Analysis (EDA), and downloading PDF reports.
    \item \textbf{Activity Diagram:} A 4-swimlane UML Activity Diagram was produced to clarify responsibilities across the Frontend (Clinician), Middleware (Flask Web Server), Backend (ML Prediction Engine), and Storage (MySQL).
\end{itemize}

\subsection{Technical Stack}

The project leverages a modern, high-performance technology stack:
\begin{itemize}
  \item \textbf{Frontend Engineering:} HTML5, modern CSS3 (with custom animations), Bootstrap 5, and Vanilla JavaScript.
  \item \textbf{Data Visualization \& Aesthetics:} Plotly.js for rendering interactive charts, combined with Three.js and Vanta.js for dynamic, hardware-accelerated cybernetic backgrounds.
  \item \textbf{Audio Feedback:} Web Audio API for synthesizing Sci-Fi diagnostic audio cues without external audio files.
  \item \textbf{Machine Learning:} TensorFlow/Keras for the CNN architecture, scikit-learn for data preprocessing, and pyts for Gramian Angular Summation Field (GASF) transformations.
  \item \textbf{Backend \& Integration (Future Work):} Flask (Python) for the RESTful API bridge between the ML engine and the Frontend SPA.
\end{itemize}

\section{Implementation Details}

\subsection{Module 1: UI Prototyping and Frontend Architecture}

The frontend Single Page Application (SPA) prototype has been successfully implemented, focusing heavily on a premium medical cybernetic aesthetic. The UI was designed to minimize cognitive load while maximizing clinical engagement.

\subsubsection{Visual Aesthetics and Theming}
We integrated glassmorphism techniques---utilizing semi-transparent backgrounds with backdrop blur filters---to provide tactile feedback on diagnostic panels. A sophisticated dark theme (\texttt{\#050811}) was chosen to reduce eye strain for clinicians. The background features an interactive 3D Vanta.js network animation that dynamically shifts colors (from neon blue to danger red) based on the computed patient risk probability. Custom typography utilizes Google Fonts \texttt{Outfit} and \texttt{Plus Jakarta Sans} for high legibility in clinical settings. Furthermore, bespoke CSS animations such as CRT scanlines and floating 3D depth (\texttt{translateZ}) elements enhance the futuristic, premium feel of the dashboard.

\subsubsection{Dynamic Input Interfaces}
The application architecture relies on three primary SPA states manipulated dynamically via JavaScript:
\begin{itemize}
    \item \textbf{Parameter Input Grid:} A responsive 3-column layout captures 13 critical physiological features (e.g., Age, Sex, Chest Pain Type, Resting BP, Cholesterol, Fasting Blood Sugar, Resting ECG, Max Heart Rate, Exercise Induced Angina, ST Depression, ST Slope, Number of Major Vessels, Thalassemia). 
    \item \textbf{Dual-Input Controls \& Clinical Ranges:} Features like Age, BP, and Cholesterol utilize synchronized range sliders and numerical inputs. Custom JavaScript logic dynamically updates "Glowing Badges" and clinical reference texts (e.g., tagging BP > 140 as "Hypertension" or Cholesterol < 200 as "Desirable") in real-time as the user adjusts the inputs.
    \item \textbf{Patient Management Mockups:} We implemented UI components for toggling between "New" and "Existing" patients, including an autocomplete datalist structure for searching patient hospital IDs.
\end{itemize}

\subsubsection{Simulated Diagnostic Pipeline}
To simulate the backend processing delay, we engineered a highly immersive "System Processing" state:
\begin{itemize}
    \item \textbf{Visual Feedback:} A non-blocking diagnostic scan overlay features complex CSS radar pulse (\texttt{radar-pulse}) and sweep animations alongside an SVG sinus rhythm drawing animation.
    \item \textbf{Scrolling Terminal:} A cascading scrolling log terminal outputs simulated computation stages (e.g., "[SYSTEM] Ingesting Telemetry Data...", "[AI] Converting 1D telemetry array to 2D Gramian Field (GASF)...").
    \item \textbf{Audio Synthesis:} We implemented a custom \texttt{SciFiAudioController} using the HTML5 Web Audio API. This synthesizes continuous sonar sweeps, data processing ticks, and a success chime upon completion, entirely through code.
\end{itemize}

\subsubsection{Results Dashboard and XAI Visualization}
The final SPA state is a comprehensive analytics module designed to present the AI's findings:
\begin{itemize}
    \item \textbf{Risk Probability:} Displays a dynamic progress bar and a CSS-animated beating heart whose pulse rate correlates with the predicted risk percentage.
    \item \textbf{Performance Metrics:} UI containers for displaying Accuracy, Precision, Recall, F1-Score, and ROC-AUC.
    \item \textbf{Plotly.js Integrations:} The dashboard includes initialized Plotly.js containers prepared to render the "XAI Consensus Profile" (comparing Global Attention, Local Grad-CAM, and Patient SHAP weights) and a detailed "2D Grad-CAM Heatmap".
    \item \textbf{AI Narrative Summary:} A dedicated section for generating human-readable clinical narratives based on the model's primary diagnostic drivers.
\end{itemize}

Client-side HTML5 validations and a robust mock JSON data generator have been configured to ensure the frontend prototype is fully operational and demonstrable prior to backend integration.

\begin{figure}[H]
    \centering
    \includegraphics[width=\textwidth,height=0.45\textheight,keepaspectratio]{diagrams/ui_dash_1.png}
    \caption{Initial Data Input Dashboard with 3D Cybernetic Background.}
    \label{fig:ui_dash_1}
\end{figure}

\begin{figure}[H]
    \centering
    \includegraphics[width=\textwidth,height=0.45\textheight,keepaspectratio]{diagrams/ui_dash_2.png}
    \caption{Active Parameter Entry utilizing synchronized sliders and glowing clinical badges.}
    \label{fig:ui_dash_2}
\end{figure}

\begin{figure}[H]
    \centering
    \includegraphics[width=\textwidth,height=0.45\textheight,keepaspectratio]{diagrams/ui_dash_3.png}
    \caption{System Processing State: Radar sweep and diagnostic overlay.}
    \label{fig:ui_dash_3}
\end{figure}

\begin{figure}[H]
    \centering
    \includegraphics[width=\textwidth,height=0.45\textheight,keepaspectratio]{diagrams/ui_dash_4.png}
    \caption{Results Dashboard: Predicted Risk Probability and Patient Metrics.}
    \label{fig:ui_dash_4}
\end{figure}

\begin{figure}[H]
    \centering
    \includegraphics[width=\textwidth,height=0.45\textheight,keepaspectratio]{diagrams/ui_dash_5.png}
    \caption{XAI Consensus Profile: Comparing Grad-CAM, SHAP, and Attention Heatmaps.}
    \label{fig:ui_dash_5}
\end{figure}

\begin{figure}[H]
    \centering
    \includegraphics[width=\textwidth,height=0.45\textheight,keepaspectratio]{diagrams/ui_dash_6.png}
    \caption{Results Dashboard: Patient Narrative Summary and PDF Report Generation.}
    \label{fig:ui_dash_6}
\end{figure}

% \subsection{Module 2: Project Scheduling and Documentation}
% 
% A comprehensive Gantt chart was designed to outline the project timeline, milestones, and task dependencies for the development of Cardio AI. The primary phases and tasks for the project lifecycle were mapped out as follows:
% \begin{enumerate}
%     \item \textbf{Requirements \& Design:} Initial proposal, system architecture drafting, and UML diagram generation.
%     \item \textbf{Model Development:} Data preprocessing, training the CNN-LSTM-Attention model, and integrating XAI techniques (Grad-CAM, SHAP).
%     \item \textbf{System Integration:} Developing the Flask REST API and setting up the MySQL database schemas.
%     \item \textbf{Frontend \& Visualization:} Building the interactive Plotly.js dashboards for explainability and developing the web interface.
% \end{enumerate}
% 
% Task dependencies were explicitly tracked to highlight critical paths. For example, it was established that model training must precede API integration, and API completion is required before the frontend XAI dashboards can be fully finalized and populated with live prediction arrays.
% 
% In parallel with scheduling, we established a robust documentation architecture. A centralized \texttt{docs} repository was configured using the \texttt{mkdocs} backend. This ensures that any commit to the \texttt{master} branch of our GitHub repository triggers a CI/CD pipeline, automating the build and deployment of all project documentation and developer journals.

\subsection{Module 2: Deep Learning Baseline Pipeline (CNN)}

To establish the core prediction engine, we implemented the foundational machine learning backend up to the Convolutional Neural Network (CNN) phase:
\begin{itemize}
    \item \textbf{Data Ingestion \& Scaling:} The UCI Heart Disease dataset was ingested and all 13 clinical features were standardized using MinMaxScaler to ensure uniform gradient propagation.
    \item \textbf{Mathematical Transformation (GASF):} Since CNNs are optimized for 2D spatial data, we transformed the 1D tabular arrays into 2D Gramian Angular Summation Field (GASF) images. This cutting-edge technique encodes temporal/feature correlations into a visual format.
    \item \textbf{CNN Architecture:} A Keras Sequential model was constructed using \texttt{Conv2D} and \texttt{MaxPooling2D} layers to extract spatial features from the GASF matrix, culminating in a \texttt{Dense} output layer with a Sigmoid activation for binary risk classification.
    \item \textbf{Initial Evaluation:} The baseline CNN was trained successfully, with loss curves and confusion matrices plotted to verify model convergence before the future integration of LSTM and Explainable AI (Grad-CAM/SHAP) components.
\end{itemize}

\section{System Modeling and Project Diagrams}

To provide a comprehensive blueprint of the Cardio AI system, several architectural and behavioral models were developed.

\subsection{Project Scheduling and Timeline (Gantt Chart)}
\begin{figure}[H]
    \centering
    \includegraphics[width=\textwidth,height=0.25\textheight,keepaspectratio]{diagrams/ganttchart.pdf}
    \caption{Project Scheduling and Timeline (Gantt Chart)}
    \label{fig:gantt}
\end{figure}

\subsection{UML Use Case Diagram}
\begin{figure}[H]
    \centering
    \includegraphics[width=\textwidth,height=0.95\textheight,keepaspectratio]{diagrams/use case diagram.pdf}
    \caption{UML Use Case Diagram identifying system boundaries and primary actors.}
    \label{fig:usecase}
\end{figure}

\subsection{UML Activity Diagram}
\begin{figure}[H]
    \centering
    \includegraphics[width=\textwidth,height=0.95\textheight,keepaspectratio]{diagrams/activity diagram.pdf}
    \caption{UML Activity Diagram mapping the procedural flow of the system.}
    \label{fig:activity}
\end{figure}

\subsection{UML Activity Diagram (Swimlanes)}
\begin{figure}[H]
    \centering
    \includegraphics[width=\textwidth,height=0.95\textheight,keepaspectratio]{diagrams/activity diagram Swimlanes.pdf}
    \caption{UML Activity Diagram mapped across specific actor swimlanes (Frontend, Middleware, Backend, DB).}
    \label{fig:activity_swimlanes}
\end{figure}

\subsection{UML Sequence Diagram}
\begin{figure}[H]
    \centering
    \includegraphics[width=\textwidth,height=0.95\textheight,keepaspectratio]{diagrams/sequence diagram.pdf}
    \caption{UML Sequence Diagram depicting the chronological interaction between system components.}
    \label{fig:sequence}
\end{figure}

\subsection{UML Collaboration Diagram}
\begin{figure}[H]
    \centering
    \includegraphics[width=\textwidth,height=0.95\textheight,keepaspectratio]{diagrams/collaboration diagram.pdf}
    \caption{UML Collaboration Diagram highlighting the structural communication links between components.}
    \label{fig:collaboration}
\end{figure}


% ===== CHAPTER 3: RESULTS & EVALUATION =====

\chapter{Progress Evaluation}

\section{Current Status and Validation}

As this is the mid-semester evaluation stage, our efforts have successfully culminated in the completion of the foundational architecture, system modeling, a highly polished frontend prototype, and the core CNN machine learning baseline.

\begin{itemize}
  \item \textbf{UML Design Validation:} The comprehensive UML models (Use Case, Activity, Sequence, and Collaboration) have been peer-reviewed and successfully validated as the logical blueprint for our distributed system.
  \item \textbf{UI/UX Testing:} The frontend components have undergone rigorous interface validation. The layout maintains responsiveness across desktop and tablet viewports, and input integrity is enforced via HTML5 attributes and JavaScript event listeners.
  \item \textbf{Mock Integration \& Flow:} The SPA transitions seamlessly between data input, processing overlays (with synchronized audio-visual feedback), and the results dashboard. The JavaScript-based mock logic successfully populates the Plotly.js charts and UI badges, proving that the frontend is fully capable of rendering the data structures we intend to send from the backend.
  \item \textbf{CNN Baseline Validation:} The foundational machine learning pipeline was successfully structured and validated in a local Python environment, proving that the clinical dataset can be correctly scaled and transformed into GASF images before being processed by the Keras CNN layers.
\end{itemize}

\section{Analysis of Progress}

Our progress strictly adheres to our initial Gantt chart objectives for the Requirements, UI, and Baseline ML phases. The robust, state-driven design of the frontend SPA ensures that integrating the Flask REST API will require minimal to no UI refactoring. 

Moving forward, the primary engineering challenge will shift to expanding the CNN baseline by adding the LSTM and Attention mechanisms to achieve our target AUC of $\geq 0.85$, and establishing the Flask bridge.


% ===== CHAPTER 4: CONCLUSION =====

\chapter{Conclusion}

\section{Summary of Work}

We have established a formidable core foundation for the Cardio AI prediction system. To date, we have:
\begin{itemize}
    \item Defined the problem scope, focusing on clinical interpretability.
    \item Architected the proposed ML pipeline and finalized the technical stack.
    \item Completed comprehensive UML modeling (Use Case, Activity, Sequence, and Collaboration diagrams).
    \item Developed a detailed project schedule using a Gantt chart.
    \item Constructed a fully interactive, glassmorphic frontend SPA prototype, complete with dynamic clinical range validations, hardware-accelerated 3D backgrounds, and synthesized audio feedback.
    \item Engineered the foundational machine learning backend up to the Convolutional Neural Network (CNN) phase, successfully integrating clinical dataset ingestion, MinMaxScaler standardization, and Gramian Angular Summation Field (GASF) image transformations.
\end{itemize}

\section{Key Findings}

During the first half of the project lifecycle, the primary engineering effort validated two core hypotheses:
\begin{itemize}
    \item \textbf{Data Transformation Efficacy:} Transforming 1D tabular clinical records into 2D Gramian Angular Summation Fields (GASF) allowed the CNN baseline to effectively extract spatial relationships among distinct clinical features, proving the viability of our imaging-based approach for non-image data.
    \item \textbf{Architectural Decoupling:} The decision to build a state-driven, heavily validated frontend SPA independently of the backend ensured rapid UI iteration. The mock integration tests prove the frontend is fully capable of rendering the complex XAI heatmaps and narrative responses that the backend will eventually serve.
\end{itemize}

\section{Future Work}

Following our established timeline, our immediate next steps encompass expanding the deep learning model and establishing backend integrations:
\begin{enumerate}
  \item \textbf{Advanced Model Training:} Expand the baseline CNN by integrating LSTM and Attention mechanisms to evaluate sequence patterns.
  \item \textbf{Explainability Engine:} Implement the core Python logic for extracting Grad-CAM heatmaps, Attention weights, and SHAP values.
  \item \textbf{System Integration:} Deploy the Flask backend to establish the RESTful bridge to our completed frontend SPA.
\end{enumerate}

\section{Final Remarks}

The Cardio AI project remains on schedule and structurally sound. By prioritizing a transparent, highly interactive user interface alongside a solid deep learning foundation, we are well-positioned to deliver an Explainable AI diagnostic tool that bridges the gap between complex neural network predictions and practical clinical trust.

% ===== BIBLIOGRAPHY =====

% Print bibliography with heading in TOC
\printbibliography[heading=bibintoc]

\end{document}
